\clearpage\corrigesChapitre{4}{Atomes, noyaux et ions}
\corrigesSouples
\begin{blocLivre}{cor-4-1}\begin{corrige}{1}{Lire une écriture nucléaire}
\begin{questions}
\item $Z=13$ protons et $A-Z=27-13=14$ neutrons.
\item 13 électrons compensent les 13 charges positives.
\end{questions}
\end{corrige}\end{blocLivre}
\begin{blocLivre}{cor-4-2}\begin{corrige}{2}{Écrire un noyau}
\begin{questions}
\item $Z=8$ et $A=8+10=18$.
\item $\ce{^{18}_{8}O}$.
\end{questions}
\end{corrige}\end{blocLivre}
\begin{blocLivre}{cor-4-3}\begin{corrige}{3}{Qualifier une entité}
\begin{questions}
\item Ne : atome ; \ce{CO2} : molécule ; \ce{Mg^{2+}} : cation ; \ce{F-} : anion.
\item \ce{Mg^{2+}}, de charge $+2e$.
\end{questions}
\end{corrige}\end{blocLivre}
\begin{blocLivre}{cor-4-4}\begin{corrige}{4}{Un anion chlorure}
\begin{questions}
\item 18 électrons : un électron a été gagné par rapport à l’atome neutre.
\item Non. L’élément reste le chlore avec 17 protons.
\end{questions}
\end{corrige}\end{blocLivre}
\begin{blocLivre}{cor-4-5}\begin{corrige}{5}{Changer d’échelle}
\schema[corrige]{atome}
\begin{questions}
\item $10^{-10}/10^{-15}=10^5$.
\item $100/10^5=\qty{0,001}{m}=\qty{1}{mm}$. Le dessin du livre n’est donc pas à l’échelle.
\end{questions}
\end{corrige}\end{blocLivre}
\begin{blocLivre}{cor-4-6}\begin{corrige}{6}{Une masse atomique}
\begin{questions}
\item $m\simeq12\times1,67\times10^{-27}=\qty{2,00e-26}{kg}$.
\item Leur masse totale vaut environ $5,47\times10^{-30}\unit{kg}$, très petite devant celle du noyau.
\end{questions}
\end{corrige}\end{blocLivre}
\begin{blocLivre}{cor-4-7}\begin{corrige}{7}{Charge du noyau}
\begin{questions}
\item $q_n=12e=\qty{1,92e-18}{C}$.
\item $q_e=-12e=\qty{-1,92e-18}{C}$. Les charges se compensent.
\end{questions}
\end{corrige}\end{blocLivre}
\begin{blocLivre}{cor-4-8}\begin{corrige}{8}{Neutraliser les charges}
\begin{questions}
\item Deux ions aluminium apportent $+6e$ et trois ions oxyde $-6e$.
\item $\ce{Al2O3}$. Les indices traduisent le rapport 2 pour 3.
\end{questions}
\end{corrige}\end{blocLivre}
\begin{blocLivre}{cor-4-9}\begin{corrige}{9}{Un ion à identifier}
\begin{questions}
\item $q=(11-10)e=+e$.
\item C’est $\ce{Na+}$ ; son noyau a $A=23$.
\end{questions}
\end{corrige}\end{blocLivre}
\begin{blocLivre}{cor-4-10}\begin{corrige}{10}{Une erreur de vocabulaire}
\begin{questions}
\item L’ion $\ce{Ca^{2+}}$ se forme par perte de deux électrons ; le noyau reste inchangé.
\item Le numéro atomique changerait : il s’agirait d’un autre élément chimique.
\end{questions}
\end{corrige}\end{blocLivre}
\begin{blocLivre}{cor-4-11}\begin{corrige}{11}{Masse et charge}
\begin{questions}
\item $q=(12-10)e=\qty{3,20e-19}{C}$.
\item Il manque deux électrons, soit environ $1,82\times10^{-30}\unit{kg}$, négligeable devant $24\times1,67\times10^{-27}\unit{kg}$. Les masses sont presque égales.
\end{questions}
\end{corrige}\end{blocLivre}
\begin{blocLivre}{cor-4-12}\begin{corrige}{12}{Un solide électriquement neutre}
\demarche{Équilibrer les charges positives et négatives, puis compter les ions.}
\begin{questions}
\item Il faut deux ions chlorure par ion calcium : $6,0\times10^{20}$ ions $\ce{Cl-}$.
\item $q=3,0\times10^{20}\times2e-6,0\times10^{20}e=0$.
\end{questions}
\retiens{La neutralité porte sur la somme des charges, pas sur une égalité des nombres d’ions.}
\end{corrige}\end{blocLivre}

